% HISTORICAL assembly fragment; the current integrated and checked text is
% parts/bmodels/bmodels.tex, which the master actually includes.
\subsection{B.4. A faithful moderated o-minimal atelier}

\begin{definition}[A precisely specified cell atelier]\label{SA21-BM-cell-definition}
Fix an o-minimal expansion \(\mathcal R\) of a real closed field \(R\),
an integer \(n\geq0\), and a definable set \(X\subseteq R^n\).
Definability always allows parameters from \(R\).
A \emph{generalized cell} is a nonempty definable set definably
homeomorphic to \((0,1)^d\) for \(d\geq0\), with \((0,1)^0\)
a singleton. No closed-ball condition is imposed on its closure.

A figure \(F\) is a finite partition of a closed definable
subset \(S\subseteq X\) into generalized cells, with frontier condition
\[
 D\cap\overline C^{\,X}\ne\varnothing
 \quad\Longrightarrow\quad D\subseteq\overline C^{\,X}
 \qquad(C,D\text{ cells of }F).
\]
The empty figure is included. Identify \(F\) with the finite
family of closed strata \(\{\overline C^{\,X}:C\in F\}\)
and write \(|F|=S\). For distinct cells \(D\subseteq\overline C^{\,X}\),
the frontier-dimension theorem
\citep[Theorem 3.22]{Coste1999} gives \(\dim D<\dim C\).
Thus their closures are distinct. The complement of \(C\)
in its closure is the finite union of the closed strata of
the lower cells, so \(C\) is relatively open and dense there.
The cell of a closed stratum \(A\) is
\[
 C_A=A\setminus\bigcup_{\substack{B\in F\\B\subsetneq A}}B.
\]
A multistrate \(m\) is the elementary figure of all strata
below one \(A\); put \(|m|=A\) and \(m^\circ=C_A\).
Thus it remembers the entire boundary decomposition.

Define \(G\leq F\) by downward-closed inclusion of closed-stratum
families, and define
\[
\begin{gathered}
 G\ll F\ \Longleftrightarrow\
 \text{each open cell of }G\text{ lies in an open cell of }F,\\
 G\preceq F\ \Longleftrightarrow\ G\ll F,\ |G|=|F|.
\end{gathered}
\]
On multistrates set \(m\perp^\circ n\) if
\(m^\circ\cap n^\circ=\varnothing\). Denote the resulting
structure by \(\mathcal A_{\mathrm{cell}}(X)\).
All closures are in the fixed ambient set \(X\), including
when \(X\) is unbounded or not closed in \(R^n\).
\end{definition}

\begin{lemma}[Proved: carriers, compatibility and restrictions]\label{SA21-BM-cell-calculus}
The relations \(\leq,\ll\) are partial orders with
\(\leq\subseteq\ll\). Refinement has a unique monotone carrier
\[
 c_{G,F}:\operatorname{Str}(G)\longrightarrow\operatorname{Str}(F),
 \qquad C_B\subseteq C_{c_{G,F}(B)}.
\]
For \(m\ll n\), the least substratum of \(n\) refined by
\(m\) is its top-cell carrier. Consequently
\[
 m\ll^\circ n\quad\Longleftrightarrow\quad
 m\ll n\text{ and }m^\circ\subseteq n^\circ.
\]
Figures are compatible exactly when, after uniting their
elementary subfigures, all distinct multistrates have disjoint
top cells; the join is then that union. They are disjoint
exactly when their supports are disjoint. For \(G\ll F\),
\(H\leq F\), induced refinement is
\[
 G|H=\{B\in G:c_{G,F}(B)\in H\},
 \qquad |G|H|=|G|\cap|H|.
\]
\end{lemma}

\begin{proof}
Reflexivity and transitivity follow from containment.
If \(F\ll G\ll F\), an \(F\)-cell \(C\) lies in a
\(G\)-cell \(D\), which lies in an \(F\)-cell \(C'\).
Nonempty intersection forces \(C=C'\), hence \(C=D\).
Repeating proves equality of figures. A downward-closed
subfamily retains its cells, so incidence implies refinement.
Carriers are unique since target cells are disjoint.
If \(B'\subseteq B\), then
\[
 C_{B'}\subseteq B'\subseteq B
 \subseteq\overline{C_{c(B)}}^{\,X}=c(B).
\]
The target cell \(C_{c(B')}\) meets \(c(B)\); frontier
compatibility gives \(c(B')\subseteq c(B)\). All carriers
of an elementary figure therefore lie below its top-cell
carrier, which must occur in any target subfigure containing
the refinement. This proves the carrier assertions.

For compatibility take the union of the proposed cells,
identifying equal elementary figures. Each original closed
stratum remains the union of its original downward family.
Frontier compatibility persists. No two distinct resulting
strata have equal support, since their disjoint top cells
would both be dense relatively open subsets of it.
Thus the union is a finite closed-support figure.
Conversely any common upper figure supplies precisely
these disjoint cells.
Disjoint supports permit their disjoint union as a figure,
with empty common incidence part. Conversely compatible
figures partition their supports by cells of their join;
a point in both lies in a common cell, and hence in a
nonempty common elementary subfigure.
Finally the inverse image of a downward family under
the carrier is downward. A point of \(|G|\) is in \(|H|\)
exactly when its carrier cell belongs to \(H\). This proves
the restriction formulas and their transitivity.
\end{proof}

\begin{theorem}[Proved: the complete axiom truth table]\label{SA21-BM-cell-axioms}
The chronological meanings of labels are fixed in Part B.1.
Positive entries hold for every definable \(X\); negative
entries mean failure for some such \(X\), proved below.
\begin{center}
\begin{tabular}{p{0.19\linewidth}p{0.57\linewidth}p{0.13\linewidth}}
\hline
System & Axioms & Value\\ \hline
156-4, Cont & Cont 1--5, Fig 1--4, finite geometric dimension,
finite principal ideals, Cond 7--8 & True\\
156-4, C & C1--C3, bounded-join primality, C5, C6,
C7\(_0\), canonical extension, C8, C8\(_0\) & True\\
156-4, C & Global incidence C4 and C4 bis & False\\
156-6 & At 1--8, At 9, At \(9'\), At 10--13,
At \(11'\), At 11 bis & True\\
156-6 & At \(9_\Sigma\) for all finite posets & True\\
156-6 & At \(9''\), or the empty/principal restriction
in \(9_\Sigma\) & False\\
156-6 & At 14, geometric At 14 bis, At 14 ter & True\\
156-6 & Literal all-target reflection in At 14 bis & False\\
156-7 & At C, At D, At D\(_0\), At CD, At CD2,
At CD2\( '\), At L1--L3, CL(spac), CL(pol), At supp & True\\
156-7 & At pol 1, At simp 1 & False\\
156-8 & Mag 1--4, strong-refinement formulation,
weaker primed clauses, Mog L1--L2, divisible-magasin clauses & True\\
156-8 & Atlas At 1--5, At Fil & True\\
156-9 & Renamed At 1--4, At spat 1--3 & True\\ \hline
\end{tabular}
\end{center}
Geometric At 14 bis means the mixed incidence condition
between a new boundary stratum and a stratum outside
that boundary. The literal all-target reading of its
uncertain transcription is treated separately below.
\end{theorem}

\begin{proof}
We verify mathematical contents, so changing labels
between chronological systems does not hide assumptions.

\emph{Multistrates.}
The orders and carriers give Mag 1--2.
Nonempty top cells make \(\perp^\circ\) symmetric and
irreflexive. Containment of top cells makes it hereditary
under strong refinement, proving Mag 3.
Distinct strata under one top stratum have disjoint
residual cells, proving the full all-pairs Mag 4.
Strong refinement is transitive since both refinement
and top-cell containment are transitive.
These facts give the corrected strong-refinement
formulation, its weakenings and renamed At 1--4.

\emph{Figures.}
The ideal below \(F\) is the downset lattice of its
finite stratum poset. Bounded joins are unions,
irreducibles are exactly elementary figures,
every figure is their join, and these irreducibles
are prime for bounded joins. The lemma proves
pairwise compatibility detection. These give
C1--C3 and At 1--3. Refinement is detected on
elementary figures, giving At 4. If \(G\leq F\),
\(m\leq F\), and \(m\ll G\), the top cell of \(m\)
is a cell of \(F\) contained in a cell of \(G\).
They are equal, so \(m\leq G\). More generally,
if compatible multistrates \(m,n\) satisfy \(m\ll n\),
put them in one common figure; the same cell
comparison gives \(m\leq n\). This is At 6 of
156-6. Least carriers give At 5, and the already
proved transitivity of strong refinement gives
At 7. Internal disjointness on figures
is defined only on elementary figures; this and
the elementary compatibility criterion give
both clauses of atlas At 5 in 156-8, scan 21.

\emph{Lieux and disjunction.}
Every singleton exists since parameters are allowed.
They are exactly the refinement-minimal nonempty
multistrates: every nonsingleton one has a proper
singleton refinement. Full and open shadows are
\(|m|\) and \(m^\circ\), respectively.
Every open shadow has a lieu; distinct lieux
are disjoint. The identity
\[
 m^\circ\cap n^\circ=\varnothing
 \quad\Longleftrightarrow\quad
 (\forall p\in m^\circ)(\forall q\in n^\circ)\ p\ne q
\]
and its whole-support analogue prove Mag L1--L2,
the displayed divisible-magasin conditions,
At 8--10, At L1--L3 and the lieu detection
clauses of the Cont systems.
Support inclusion makes disjointness hereditary
under refinement.

\emph{Compatibility and gluing.}
For figures \(F,G\), intersect their sets of elementary
subfigures to obtain their greatest common incidence
subfigure \(K\). Then
\[
 F\bowtie G\quad\Longleftrightarrow\quad
 (|F|\setminus|K|)\cap(|G|\setminus|K|)=\varnothing.
\]
The forward direction uses their common partition.
For the reverse, outside cells are disjoint by
the formula, inside cells are identical, and
outside cells cannot meet an inside cell.
The union criterion applies. This proves
CL(spac), At \(9'\), and At \(9_\Sigma\)
with all finite common posets allowed.

Suppose \(F\bowtie G\), \(F'\ll F\), \(G'\ll G\),
and \(F'|K\bowtie G'|K\). Any intersection
of a refining \(F'\)-cell and \(G'\)-cell
lies in \(|F|\cap|G|=|K|\). A refining cell
meeting \(|K|\) has its carrier in \(K\),
so belongs to the appropriate restriction.
All possible intersections are thus governed
by compatible restrictions. The union
criterion proves \(F'\bowtie G'\):
this is At 12 and At C.
The same restriction identity with
empty intersection proves At D and D\(_0\).
For absolute At CD2, suppose all refinements
of \(F,G\) whose restrictions to \(K\) are empty
are pairwise disjoint. Singleton points outside
\(|K|\) are among these refinements. Hence
\(|F|\setminus|K|\) and \(|G|\setminus|K|\)
are disjoint, and the displayed formula
gives \(F\bowtie G\). This proves At CD2;
the combined At CD formulation follows
from it and the relative criterion.
CL(pol) and At CD2\( '\) add a hypothesis
that \(K\) be empty or elementary to a
sufficient criterion, so follow from
CL(spac) and At CD2. They differ from
the biconditional At \(9''\).

Refinements agreeing on overlaps glue
by union, with precisely the prescribed
restrictions, proving At 11 and At \(11'\).
For subdivisions the supports cover the
original support, proving C8.
Any incidence cover of a finite figure
has a finite subcover. The data on
any further cover member agree with
the finite gluing on its intersections
with this subcover, which cover it.
Thus the arbitrary-cover version holds.
For At 11 bis, a directed family
with an existing incidence supremum
on one stratum is stationary there:
the supremum has finitely many
strata, included at one directed index.
Finitely many original strata give
one global stabilization index.
The eventual figure is the required
supremum. Starting instead with
stratumwise stationarity gives At Fil.

\emph{Supports and induction.}
If \(F'\ll F\) and some figure is disjoint
from \(F'\) but not \(F\), choose
\(p\in|F|\setminus|F'|\). Its singleton
refines \(F\), is disjoint from \(F'\),
and is not disjoint from \(F\).
This proves At supp, At 13 and
Cond 7 in their exterior-lieu forms.
It also proves that, for a refinement,
equal annihilators are equivalent to
equal supports, giving C5.
The restriction formula gives C6,
including uniqueness and commutation
with bounded joins. A refinement
missing a target point is strictly
enlarged by adjoining that singleton;
with equal supports no enlargement
is possible. Hence subdivisions are
exactly incidence-maximal refinements.
Antisymmetry gives C8\(_0\).
Extensions are proved separately.

\emph{Cont and dimension.}
The empty figure and points exist.
Admissibility, downward subfamilies,
compatibility detection and finite
compatible unions were proved in
the lemma. These verify the
preliminary and final Cont clauses.
Principal stratum ideals are finite.
Each strict frontier inclusion
lowers dimension by
\citet[Theorem 3.22]{Coste1999}.
For \(m'\ll m\), every chain
below \(m'\) therefore has length
at most \(\dim|m|\leq n\), since
\(|m'|\subseteq|m|\). This bounds
geometric dimension uniformly over
all refinements, as Cont 3 requires.
Empty figures, one-stratum figures
and dimension zero are included.
We do not assert that an unbounded
infinite union is a figure.
Spatiality and negative entries
are proved below.
\end{proof}
